but such that
@@M0@@

\sgaStableTarget{sga6.e14.para.5-3}\textbf{5.3.} When @@M1@@ is smooth and quasi-projective over a field @@M2@@, one can likewise ask about the existence of a ring homomorphism
@@M3@@
with values in the Chow ring @@M4@@, taking Chern classes to Chern classes. One can repeat the preceding considerations, which show in particular how to construct @@M5@@ directly in the absence of @@M6@@ itself.

\sgaStableTarget{sga6.e14.para.5-4}\textbf{5.4.} To conclude this section, observe that essentially the same question arises when, in place of a scheme @@M7@@, one considers a compact topological space @@M8@@, say; one can again consider the associated graded @@M9@@ of @@M10@@ endowed with its @@M11@@-filtration, and seek a ring homomorphism
@@M12@@
taking Chern classes to Chern classes.

\sgaStableTarget{sga6.e14.sec.6}\section*{6. Cohomological Riemann--Roch theorem and Gysin homomorphism}

\sgaStableTarget{sga6.e14.para.6-1}\textbf{6.1.} For simplicity, we restrict ourselves here to schemes @@M13@@ admitting an ample invertible Module; this implies that the naive @@M14@@ and the other one can be identified. Let us begin with the Riemann--Roch formula ``with denominators'' for a proper complete-intersection morphism @@M15@@. Since (for a prime @@M16@@ prime to the residue characteristics of @@M17@@) one can define the @@M18@@-adic Chern classes of a locally free Module @@M19@@,
@@M20@@
(SGA 5 VII), corresponding to a homomorphism
\sgaStableTarget{sga6.e14.eq.6-1}@@M21@@
the question arises of giving a cohomological variant of the Riemann--Roch formula. The only ingredient missing to give the formula a meaning is a \uline{Gysin homomorphism}
